%------------------------------------------------------------------------------%

\usepackage{integracao_e_series}

\title{Séries de Potências -- Introdução}

%------------------------------------------------------------------------------%
\begin{document}

\addtitle

%------------------------------------------------------------------------------%
\begin{frame}{Séries de Funções}
  Queremos estudar séries da forma
  \[
    F(x) = \sum_{n=0}^\infty\; f_n(x)
  \]
  \pause
  onde $f_n:D \subset \R\to\R$ são \emph{funções} da variável real $x$ \\[5mm]

  \pause
  Para cada $x$ fixo temos uma série numérica \\[5mm]

  \pause
  Queremos saber para quais valores de $x$ podemos somar a série \\[5mm]

  \pause
  O domínio de $F$ é o conjunto de valores de $x$ para os quais a série converge
\end{frame}

%------------------------------------------------------------------------------%
\begin{frame}{Séries de Funções}
  \begin{description}[<+->]
    \item[Séries de Potências] Séries de potências de $x$
    \[
      \sum_{n=0}^\infty a_n x^n
    \]

    \item[Séries de Fourier] Série de funções trigonométricas
    \[
      \sum_{n=0}^\infty a_n \cos(nx) + b_n \sin(nx)
    \]
  \end{description}
\end{frame}

%------------------------------------------------------------------------------%
\begin{frame}{Objetivo}
  Vamos mostrar que
  \pause
  \[
    \sin(x) = x - \frac{x^3}{3!} + \frac{x^5}{5!} + \frac{x^7}{7!} - \cdots
  \]
  \pause
  \[
    \cos(x) = 1 - \frac{x^2}{2!} + \frac{x^4}{4!} + \frac{x^8}{8!} - \cdots
  \]
  \pause
  \[
    e^x = 1 + x + \frac{x^2}{2!} + \frac{x^3}{3!} + \frac{x^4}{4!} +  \cdots
  \]
\end{frame}

%------------------------------------------------------------------------------%
\begin{frame}{Identidade de Euler}
  \Huge
  \[
    e^{\pi i} + 1 = 0
  \]
\end{frame}

%------------------------------------------------------------------------------%
\end{document}
%------------------------------------------------------------------------------%
